\documentclass[a4paper,11pt]{article}

\usepackage[czech]{babel}
\usepackage[T1]{fontenc}
\usepackage[utf8]{inputenc}
\usepackage{amsmath,amssymb}
\usepackage{geometry}
\usepackage{enumitem}

\geometry{a4paper,margin=17mm}
\pagestyle{empty}
\setlength{\parindent}{0pt}
\setlength{\parskip}{4pt}
\setlist[enumerate]{leftmargin=*,itemsep=4pt,topsep=4pt}

\newcommand{\TestTitle}{Komplexní čísla}
\newcommand{\TestHeader}[1]{%
    \begin{center}
        {\Large\bfseries TEST -- \TestTitle}\\[2mm]
        {\bfseries Varianta #1}
    \end{center}
    \vspace{2mm}
    Jméno a příjmení: \rule{9cm}{0.4pt}\hfill Třída: \rule{2cm}{0.4pt}\\[3mm]
    Datum: \rule{3cm}{0.4pt}\hfill Počet bodů: \rule{2cm}{0.4pt}\hfill Známka: \rule{2cm}{0.4pt}\\[2mm]
    \hrule\vspace{4mm}
}
\newcommand{\AnswerSpace}[1]{\vspace{#1}}

\begin{document}
\TestHeader{G}

% ZADÁNÍ VARIANTY G
\begin{enumerate}
    \item Řešte v $\mathbb C$: $z\overline z-z=6-2i$.
    \AnswerSpace{20mm}
    \item Určete všechna komplexní čísla $z$, pro která platí
    $|z-(1+i)|=|z-(3-i)|$ a $|z|=\sqrt{10}$.
    Jejich obrazy znázorněte v Gaussově rovině.
    \AnswerSpace{38mm}
    \item Řešte v $\mathbb C$: $x^2+(2-i)x+3-i=0$.
    \AnswerSpace{23mm}
    \item Bez postupného roznásobování vypočtěte
    $\displaystyle\left(\frac{1+\sqrt3i}{1+i}\right)^6$.
    Výsledek zapište v algebraickém tvaru.
    \AnswerSpace{26mm}
    \item Řešte v $\mathbb C$: $z^3-64i=0$.
    Všechny kořeny zapište v goniometrickém i algebraickém tvaru
    a znázorněte je v Gaussově rovině.
    \AnswerSpace{44mm}
\end{enumerate}

\newpage
\TestHeader{H}

% ZADÁNÍ VARIANTY H
\begin{enumerate}
    \item Řešte v $\mathbb C$: $z\overline z+z=6+2i$.
    \AnswerSpace{20mm}
    \item Určete všechna komplexní čísla $z$, pro která platí
    $|z-(1-i)|=|z-(3+i)|$ a $|z|=\sqrt{10}$.
    Jejich obrazy znázorněte v Gaussově rovině.
    \AnswerSpace{38mm}
    \item Řešte v $\mathbb C$: $ix^2-3x+4i=0$.
    \AnswerSpace{23mm}
    \item Bez postupného roznásobování vypočtěte
    $\displaystyle\left(\frac{1-\sqrt3i}{1-i}\right)^6$.
    Výsledek zapište v algebraickém tvaru.
    \AnswerSpace{26mm}
    \item Řešte v $\mathbb C$: $z^3+8i=0$.
    Všechny kořeny zapište v goniometrickém i algebraickém tvaru
    a znázorněte je v Gaussově rovině.
    \AnswerSpace{44mm}
\end{enumerate}
\end{document}
